\documentclass[../../../hsm_tok_q.tex]{subfiles}

\begin{document}
    \setcounter{figure}{0}
    Kさんのクラスでは，先生が示した問題をみんなで考えた．

    次の各問に答えよ．

    \begin{tcolorbox}[title=［先生が示した問題］,
        colback=white, coltitle=black, colbacktitle=white,colframe=black, 
        toptitle=1pt, bottomtitle=1pt,
        arc=4pt, boxrule=0.5pt, before upper={\parindent=1\zw},
        left=3pt, right=3pt, top=3pt, bottom=4pt]

        \textsf{図\ref{fig/hsm_tok_2022_2nd_02_01_q}}に示した立体は，%
        底面が半径2\,cmの円，高さが5\,cmの円柱で，底面の2つの円の中心をそれぞれO，Pとし，%
        点Oと点Pを結んでできる線分は2つの底面に垂直である．
       
        円Oの周上にある点をQとし，点Oと点P，点Oと点Qをそれぞれ結び，点Qを通り線分OPに平行な直線を引き，%
        円Pの周上で交わる点をRとし，点Pと点Rを結ぶ．
        
        四角形OPRQの面積を$\mathrm{S\,cm^2}$，円柱の体積を$\mathrm{V\,cm^3}$とするとき，%
        VをSを用いた式で表しなさい．
        \\
        \begin{minipage}{\linewidth}
            \vspace{.5\baselineskip}
            {\centering
                \setlength{\abovecaptionskip}{5pt}
                \includegraphics%
                    {\subfix{fig_hsm_tok_2022_2nd_02/fig_hsm_tok_2022_2nd_02_01_q.pdf}}
                    \captionof{figure}{} \label{fig/hsm_tok_2022_2nd_02_01_q}
            }
        \end{minipage}
    \end{tcolorbox}
    %
    Kさんは，［先生が示した問題］の答えを次の形の式で表した．%
    Kさんの答えは正しかった．

    \vspace{.4\baselineskip}
    〈Kさんの答え〉 V＝\:\myboxa{　　　}\:S
    
    \vspace{.4\baselineskip}
    \begin{enumerate}
        \item 〈Kさんの答え〉の\:\myboxa{　　　}\:に当てはまるものを，%
            次の{\textsf ア}〜\textsf{エ}のうちから選び，記号で答えよ．

            ただし，円周率は$\pi$とする．
    
            \begin{tabularx}{\dimexpr\linewidth-\zw}{@{}LLLL@{}}
                {\textsf ア}\quad$2\pi$ & {\textsf イ}\quad$4\pi$%
                & {\textsf ウ}\quad$10\pi$ & {\textsf エ}\quad$20\pi$
            \end{tabularx}
    \end{enumerate}

    \newpage
    
    Kさんのグループは，［先生が示した問題］をもとにして，次の問題を考えた．

    \begin{tcolorbox}[title=［Kさんのグループが作った問題］,
        colback=white, coltitle=black, colbacktitle=white,colframe=black, 
        toptitle=1pt, bottomtitle=1pt,
        arc=4pt, boxrule=0.5pt, before upper={\parindent=1\zw},
        left=3pt, right=3pt, top=3pt, bottom=4pt]

        $a$，$h$を正の数とする．

        \textsf{図\ref{fig/hsm_tok_2022_2nd_02_02_q}}に示した立体ABCD--EFGHは，%
        底面が1辺$2a\,\mathrm{cm}$の正方形，高さが$h\,\mathrm{cm}$の直方体である．

        四角形ABCDにおいて，対角線ACと対角線BDとの交点をO，四角形EFGHにおいて，
        対角線EGと対角線FHとの交点をPとする．

        辺BC，辺FGの中点をそれぞれQ，Rとし，点Oと点P，点Oと点Q，点Pと点R，点Qと点Rをそれぞれ結ぶ．

        四角形ABCDの周の長さを$\ell\,\mathrm{cm}$，四角形OPRQの面積を$\mathrm{S\,cm^2}$，%
        立体ABCD--EFGHの体積を$\mathrm{V\,cm^3}$とするとき，%
        $\mathrm{V}=\myfraca{1}{2}\ell\,\mathrm{S}$となることを確かめてみよう．\rule{0pt}{16pt}
        \\
        \begin{minipage}{\linewidth}
            \vspace{.5\baselineskip}
            {\centering

            \setlength{\abovecaptionskip}{2pt}
                \includegraphics%
                    {\subfix{fig_hsm_tok_2022_2nd_02/fig_hsm_tok_2022_2nd_02_02_q.pdf}}
                    \captionof{figure}{} \label{fig/hsm_tok_2022_2nd_02_02_q}
            }
        \end{minipage}
        %
    \end{tcolorbox}
    %
    \begin{enumerate}[resume]
        \item ［Kさんのグループが作った問題］で，%
            $\mathrm{V}=\myfraca{1}{2}\ell\,\mathrm{S}$となることを証明せよ．
    \end{enumerate}
\end{document}

